\documentclass[11pt]{article}
% Shared preamble for manuscripts-v2 papers. Load with % Shared preamble for manuscripts-v2 papers. Load with % Shared preamble for manuscripts-v2 papers. Load with \input{../common/preamble.tex}
% immediately after \documentclass[11pt]{article}.
\ifdefined\pdftrailerid\pdftrailerid{}\fi
\usepackage[T1]{fontenc}
\usepackage{lmodern}
\usepackage[letterpaper,margin=1.02in]{geometry}
\usepackage{amsmath,amssymb,amsthm,mathtools,booktabs,array,microtype,enumitem}
\usepackage{tikz}
\usepackage{xurl}
\usepackage[colorlinks=true,allcolors=bluegreen]{hyperref}
\definecolor{bluegreen}{RGB}{22,83,91}
\usepackage[capitalize,nameinlink]{cleveref}

\theoremstyle{plain}
\newtheorem{theorem}{Theorem}[section]
\newtheorem{lemma}[theorem]{Lemma}
\newtheorem{proposition}[theorem]{Proposition}
\newtheorem{corollary}[theorem]{Corollary}
\newtheorem{conjecture}[theorem]{Conjecture}
\newtheorem*{theorem*}{Theorem}
\newtheorem{maintheorem}{Theorem}
\renewcommand{\themaintheorem}{\Alph{maintheorem}}
\theoremstyle{definition}
\newtheorem{definition}[theorem]{Definition}
\newtheorem{example}[theorem]{Example}
\newtheorem{question}[theorem]{Question}
\theoremstyle{remark}
\newtheorem{remark}[theorem]{Remark}
\crefname{maintheorem}{Theorem}{Theorems}
\crefname{conjecture}{Conjecture}{Conjectures}
\crefname{question}{Question}{Questions}
\crefname{example}{Example}{Examples}

\newcommand{\R}{\mathbb R}
\newcommand{\Q}{\mathbb Q}
\newcommand{\Z}{\mathbb Z}
\newcommand{\N}{\mathbb N}
\newcommand{\T}{\mathbb T}
\newcommand{\C}{\mathbb C}
\DeclareMathOperator{\den}{den}
\newcommand{\e}{\mathrm e}

% Immutable Lean source links at research-checkpoint-18.
\newcommand{\leancommit}{e4809cf561573aef588e6ffc491a35cdb79fd9ec}
\newcommand{\leanbase}{https://github.com/eldernos/MathematicsResearch/blob/\leancommit}
% \lean{File}{line}{label}: link to NormalityChecker/File.lean#Lline at the pinned commit.
\newcommand{\lean}[3]{\href{\leanbase/NormalityChecker/#1.lean\#L#2}{#3}}
% \repo{path}{label}: link to any repository path at the pinned commit.
\newcommand{\repo}[2]{\href{https://github.com/eldernos/MathematicsResearch/tree/\leancommit/#1}{#2}}

% Evidence labels (see papers/WAVE2-CONTRACT.md).
\newcommand{\evidence}[1]{\textup{[\emph{#1}]}}

\setlength{\emergencystretch}{1.5em}
\raggedbottom

% immediately after \documentclass[11pt]{article}.
\ifdefined\pdftrailerid\pdftrailerid{}\fi
\usepackage[T1]{fontenc}
\usepackage{lmodern}
\usepackage[letterpaper,margin=1.02in]{geometry}
\usepackage{amsmath,amssymb,amsthm,mathtools,booktabs,array,microtype,enumitem}
\usepackage{tikz}
\usepackage{xurl}
\usepackage[colorlinks=true,allcolors=bluegreen]{hyperref}
\definecolor{bluegreen}{RGB}{22,83,91}
\usepackage[capitalize,nameinlink]{cleveref}

\theoremstyle{plain}
\newtheorem{theorem}{Theorem}[section]
\newtheorem{lemma}[theorem]{Lemma}
\newtheorem{proposition}[theorem]{Proposition}
\newtheorem{corollary}[theorem]{Corollary}
\newtheorem{conjecture}[theorem]{Conjecture}
\newtheorem*{theorem*}{Theorem}
\newtheorem{maintheorem}{Theorem}
\renewcommand{\themaintheorem}{\Alph{maintheorem}}
\theoremstyle{definition}
\newtheorem{definition}[theorem]{Definition}
\newtheorem{example}[theorem]{Example}
\newtheorem{question}[theorem]{Question}
\theoremstyle{remark}
\newtheorem{remark}[theorem]{Remark}
\crefname{maintheorem}{Theorem}{Theorems}
\crefname{conjecture}{Conjecture}{Conjectures}
\crefname{question}{Question}{Questions}
\crefname{example}{Example}{Examples}

\newcommand{\R}{\mathbb R}
\newcommand{\Q}{\mathbb Q}
\newcommand{\Z}{\mathbb Z}
\newcommand{\N}{\mathbb N}
\newcommand{\T}{\mathbb T}
\newcommand{\C}{\mathbb C}
\DeclareMathOperator{\den}{den}
\newcommand{\e}{\mathrm e}

% Immutable Lean source links at research-checkpoint-18.
\newcommand{\leancommit}{e4809cf561573aef588e6ffc491a35cdb79fd9ec}
\newcommand{\leanbase}{https://github.com/eldernos/MathematicsResearch/blob/\leancommit}
% \lean{File}{line}{label}: link to NormalityChecker/File.lean#Lline at the pinned commit.
\newcommand{\lean}[3]{\href{\leanbase/NormalityChecker/#1.lean\#L#2}{#3}}
% \repo{path}{label}: link to any repository path at the pinned commit.
\newcommand{\repo}[2]{\href{https://github.com/eldernos/MathematicsResearch/tree/\leancommit/#1}{#2}}

% Evidence labels (see papers/WAVE2-CONTRACT.md).
\newcommand{\evidence}[1]{\textup{[\emph{#1}]}}

\setlength{\emergencystretch}{1.5em}
\raggedbottom

% immediately after \documentclass[11pt]{article}.
\ifdefined\pdftrailerid\pdftrailerid{}\fi
\usepackage[T1]{fontenc}
\usepackage{lmodern}
\usepackage[letterpaper,margin=1.02in]{geometry}
\usepackage{amsmath,amssymb,amsthm,mathtools,booktabs,array,microtype,enumitem}
\usepackage{tikz}
\usepackage{xurl}
\usepackage[colorlinks=true,allcolors=bluegreen]{hyperref}
\definecolor{bluegreen}{RGB}{22,83,91}
\usepackage[capitalize,nameinlink]{cleveref}

\theoremstyle{plain}
\newtheorem{theorem}{Theorem}[section]
\newtheorem{lemma}[theorem]{Lemma}
\newtheorem{proposition}[theorem]{Proposition}
\newtheorem{corollary}[theorem]{Corollary}
\newtheorem{conjecture}[theorem]{Conjecture}
\newtheorem*{theorem*}{Theorem}
\newtheorem{maintheorem}{Theorem}
\renewcommand{\themaintheorem}{\Alph{maintheorem}}
\theoremstyle{definition}
\newtheorem{definition}[theorem]{Definition}
\newtheorem{example}[theorem]{Example}
\newtheorem{question}[theorem]{Question}
\theoremstyle{remark}
\newtheorem{remark}[theorem]{Remark}
\crefname{maintheorem}{Theorem}{Theorems}
\crefname{conjecture}{Conjecture}{Conjectures}
\crefname{question}{Question}{Questions}
\crefname{example}{Example}{Examples}

\newcommand{\R}{\mathbb R}
\newcommand{\Q}{\mathbb Q}
\newcommand{\Z}{\mathbb Z}
\newcommand{\N}{\mathbb N}
\newcommand{\T}{\mathbb T}
\newcommand{\C}{\mathbb C}
\DeclareMathOperator{\den}{den}
\newcommand{\e}{\mathrm e}

% Immutable Lean source links at research-checkpoint-18.
\newcommand{\leancommit}{e4809cf561573aef588e6ffc491a35cdb79fd9ec}
\newcommand{\leanbase}{https://github.com/eldernos/MathematicsResearch/blob/\leancommit}
% \lean{File}{line}{label}: link to NormalityChecker/File.lean#Lline at the pinned commit.
\newcommand{\lean}[3]{\href{\leanbase/NormalityChecker/#1.lean\#L#2}{#3}}
% \repo{path}{label}: link to any repository path at the pinned commit.
\newcommand{\repo}[2]{\href{https://github.com/eldernos/MathematicsResearch/tree/\leancommit/#1}{#2}}

% Evidence labels (see papers/WAVE2-CONTRACT.md).
\newcommand{\evidence}[1]{\textup{[\emph{#1}]}}

\setlength{\emergencystretch}{1.5em}
\raggedbottom

% Author information approved by the repository owner on 2026-11-25:
%   name: Christopher Anderson; designation: Independent Researcher.
% No academic affiliation, ORCID or personal contact details have been supplied.
% Do not add any that the owner has not approved.
\newcommand{\OwnerAuthor}{Christopher Anderson}
\newcommand{\OwnerAffiliation}{Independent Researcher}
\newcommand{\OwnerContact}{Corrections and comments: \url{https://github.com/eldernos/MathematicsResearch/issues}}
\newcommand{\ReleaseLabel}{manuscripts-v2.0.0}
\newcommand{\LicenseLine}{\copyright\ 2026 Christopher Anderson. Licensed under \href{https://creativecommons.org/licenses/by/4.0/}{CC BY 4.0}.}
% \makebyline{version}{date}: centred author block used under the title.
\newcommand{\makebyline}[2]{%
  {\large \OwnerAuthor\par}%
  \ifx\OwnerAffiliation\empty\else{\small \OwnerAffiliation\par}\fi
  \vspace{0.3em}
  {\small Version #1 \enspace$\cdot$\enspace #2 \enspace$\cdot$\enspace \ReleaseLabel\par}%
  {\footnotesize \OwnerContact\par}%
  {\footnotesize \LicenseLine\par}%
  \hypersetup{pdfauthor={\OwnerAuthor}}}

\hypersetup{pdftitle={The Arithmetic of Extreme Discrepancy: Translation, Addition Spectra and Rational Operations}}
\newcommand{\dist}[1]{\lVert #1\rVert}
\newcommand{\fr}[1]{\{#1\}}
\DeclareMathOperator{\ord}{ord}
\newcommand{\ttt}[1]{\texttt{#1}}
\newcommand{\leanP}[3]{\leanat{\phaseXIXcommit}{#1}{#2}{#3}}
\setlength{\emergencystretch}{3em}

\begin{document}
\begin{center}
{\LARGE\bfseries The Arithmetic of Extreme Discrepancy:\\[3pt] Translation, Addition Spectra\\[3pt] and Rational Operations\par}
\vspace{10pt}
\makebyline{1.0}{25 November 2026}
\end{center}

\begin{abstract}
For an integer base $b\ge2$ let $D_N(x)$ be the extreme discrepancy of $x,bx,\ldots,b^{N-1}x$ modulo
one. We prove the finite identity $\sup_x|D_N(x+z)-D_N(x)|=1-D_N(z)$, attained at $x=0$, and
classify the translations that preserve both limiting discrepancies of every real number: they
are the numbers $c+w$ with $(b-1)c\in\Z$ and $N^{-1}\sum_{n<N}\dist{b^nw}\to0$. The limits of
$D_N(x)$, $D_N(y)$ and $D_N(x+y)$ are exactly the triples in $[0,1]^3$ with
$\max(0,\alpha+\beta-1)\le\gamma\le1-|\alpha-\beta|$. Every such triple is realized by positive
rational expansions on a cubic schedule with simultaneous error $O(N^{-1/3})$, uniformly in the
triple, by computable and algebraically independent witnesses; the analogous deficit polytope for
$d$ summands is determined. The translation deficit $1-D_N(x-y)$ is a pseudometric, an exact form
of L\'evy's concentration function comparable to the Ky Fan distance to constants. The rational
scalars preserving its asymptotic kernel are exactly $\Z[1/b]$, an individual rational scaling is
decided by a residue-density test, and on the boundary $D_N(mx)\to1$ the attainable limits of
$D_N(x)$ form $[1/m_0,1]$, where $m_0$ is the part of $m$ prime to $b$. The finite identity, the
pseudometric axioms and the qualitative addition spectrum are verified by the Lean~4 kernel.
\end{abstract}

\section{Introduction}\label{sec:intro}

Throughout, $b\ge2$ is an integer, $\T=\R/\Z$, $\fr{y}$ is the fractional part of $y$ and
$\dist{y}$ its distance to the nearest integer. For real $x$ and $N\ge1$ the \emph{extreme
discrepancy} of the first $N$ points of the radix orbit of $x$ is
\[
D_N(x)=D_N(x;b)=\sup_{0\le a<c\le1}\left|\frac{\#\{0\le n<N:\fr{b^nx}\in[a,c)\}}{N}-(c-a)\right|.
\]
A real number is normal to base $b$ if and only if $D_N(x)\to0$ \cite[Chapter~1, Theorem~8.1,
p.~70, and Chapter~2, Theorem~1.1, p.~89]{KN}. At the other extreme, $D_N(x)=1$ for every $N$ when
the orbit of $x$ is constant. This paper studies how $D_N$ behaves under the arithmetic operations
of $\R$: translation, addition of two inputs, and multiplication by rationals.

\begin{maintheorem}[Sharp translation identity]\label{thm:A}
For every real $z$ and every $N\ge1$,
\[
\sup_{x\in\R}\bigl|D_N(x+z)-D_N(x)\bigr|=1-D_N(z),
\]
and the supremum is attained at $x=0$. Moreover the following are equivalent: $D_N(z)\to1$;
$D_N(x+z)-D_N(x)\to0$ for every real $x$; translation by $z$ preserves the lower and the upper
limit of $D_N(x)$ for every real $x$; and $z=c+w$ with $(b-1)c\in\Z$ and
$N^{-1}\sum_{n<N}\dist{b^nw}\to0$.
\end{maintheorem}

The identity is a Lean declaration checked by the kernel (\Cref{thm:sharp}); the classification is
\Cref{thm:preserve}. Rauzy characterized the numbers $z$ for which $x+z$ is normal whenever $x$ is
normal \cite{Rauzy76}; Bergelson and Downarowicz identify them with the deterministic numbers
\cite[Definitions~3.5 and~3.10, Theorem~3.12]{BD25}. \Cref{thm:A} is the counterpart for numerical
discrepancy, and its class is strictly smaller (\Cref{rem:rauzy}).

Two applications of \Cref{thm:A} constrain the discrepancies of a sum. Nothing else does.

\begin{maintheorem}[Addition spectrum]\label{thm:B}
For all real $x,y$ and every $N\ge1$, the numbers $\alpha=D_N(x)$, $\beta=D_N(y)$ and
$\gamma=D_N(x+y)$ satisfy
\begin{equation}\label{eq:facets}
\max(0,\alpha+\beta-1)\le\gamma\le1-|\alpha-\beta| .
\end{equation}
Conversely every $(\alpha,\beta,\gamma)\in[0,1]^3$ satisfying \eqref{eq:facets} is the triple of
limits of $D_N(x)$, $D_N(y)$, $D_N(x+y)$ for some real $x,y$. For every odd prime $p\nmid b$ with
$t=\ord_p(b)$ there is $C_{b,p}$ such that each such triple is realized by
$x=\sum_kc_kb^{-tk^3}$ and $y=\sum_kd_kb^{-tk^3}$ with rationals $c_k,d_k\in(0,1]$ of $p$-power
denominator and simultaneous error at most $C_{b,p}N^{-1/3}$ for every $N\ge1$. The witnesses can
be chosen algebraically independent over $\Q$, with the same rate, and they are computable
uniformly from names of a computable triple.
\end{maintheorem}

Necessity at every $N$ and the qualitative sufficiency in every base are kernel-verified
(\Cref{thm:addition}). The rate, the computability and the algebraic independence are written
proofs that use one kernel-verified lemma (\Cref{thm:additionrate,thm:independent}). For $d$
summands the deficits $1-\delta$ of the $d+1$ limits are exactly the points of the unit cube in which
each coordinate is at most the sum of the others (\Cref{thm:dsummands}).

The deficit $1-D_N$ is a distance.

\begin{maintheorem}[Discrepancy distance and rational scalars]\label{thm:C}
For $N\ge1$, $d_N(x,y)=1-D_N(x-y)$ is a translation-invariant pseudometric on $\R$, equal to
$\max_w|D_N(w+x-y)-D_N(w)|$; for $N\ge2$ it vanishes exactly when $(b-1)(x-y)\in\Z$. The set
$K_b=\{z:D_N(z)\to1\}$ is a subgroup of $\R$, the kernel of $\limsup_Nd_N$. For a reduced fraction
$a/d$ one has $(a/d)K_b\subseteq K_b$ if and only if every prime factor of $d$ divides $b$. Whether
$(a/d)z\in K_b$ for an individual $z\in K_b$ is decided by the density of a single residue class
modulo $d$ of the integers nearest to $b^nz-c$, where $c$ is the limiting fixed point of the orbit.
\end{maintheorem}

The pseudometric axioms and the zero set are kernel-verified (\Cref{prop:metric}); the scalar
results are \Cref{thm:scalars,prop:residue}. Multiplication by integers is controlled as follows.

\begin{maintheorem}[Integer dilation]\label{thm:D}
Let $m\ne0$ be an integer. For every real $x$ and $N\ge1$,
$\max(0,1-|m|(1-D_N(x)))\le D_N(mx)\le\min(1,|m|D_N(x))$. Write $|m|=m_0m_b$ with
$\gcd(m_0,b)=1$ and every prime factor of $m_b$ dividing $b$. If $D_N(mx)\to1$, then
$\liminf_ND_N(x)\ge1/m_0$; conversely every $\delta\in[1/m_0,1]$ is $\lim_ND_N(x)$ for some $x$ with
$D_N(mx)\to1$, realized by positive rational expansions on a cubic schedule at rate $O(N^{-1/3})$.
\end{maintheorem}

These are \Cref{prop:dilation,thm:boundary}; the lower bound in the first statement is
kernel-verified. By contrast, nonzero rational multiplication preserves normality, and even the lower
and upper point entropies, exactly \cite[Proposition~4.25]{BD25}; this refines a classical theorem of
Wall.

\paragraph{Relation to earlier work and scope of the contribution.}
The objects in \Cref{thm:A,thm:C} are classical. The extreme discrepancy on $[0,1)$ is invariant
under rotation of the circle and equals Kuiper's circular statistic $V_N=D_N^++D_N^-$
\cite{Kuiper60}, which Kuiper introduced because the Kolmogorov--Smirnov statistic is not invariant
under rotations of the circle. The quantity $1-D$ is an exact rewriting of L\'evy's
concentration function on the circle \cite{Levy39}: if $Q_\mu(l)$ is the largest mass of an arc of
length $l$, then $1-D(\mu)=\inf_l(1-Q_\mu(l)+l)$. In probabilistic terms $1-D(\mu)$ lies between the
Ky Fan distance \cite{Fan43} from a random point with law $\mu$ to the nearest constant and three
times that distance, and the Ky Fan distance to a constant $c$ is, up to boundary conventions, the
Lévy--Prokhorov distance from $\mu$ to $\delta_c$ \cite{Prokhorov56,Dudley02}
(\Cref{prop:concentration}). The pseudometric $d_N$ is therefore not a
new probability metric. What we add, in the sources we examined, is the exact finite identity of
\Cref{thm:A} for paired radix orbits, the explicit arithmetic realizations, the quantitative rates
and the formal verification.

For addition the closest precedent is the theory of point entropy. Bergelson and Downarowicz prove
that the lower and upper point entropies $\underline h,\overline h$ for multiplication by $b$ on
$\T$ satisfy, for all $x,y$, a family of inequalities that reduce, when $h(x)$ and $h(y)$ exist, to
$|h(x)-h(y)|\le\underline h(x+y)\le\overline h(x+y)\le\min(\log b,h(x)+h(y))$
\cite[Proposition~4.9 and Corollary~4.10(1)]{BD25}; the subadditivity of upper entropy is attributed
there to Kamae, and the attainability of the bounds is studied in \cite[Section~4.3]{BD25}. In the
normalized deficit coordinates $1-D$ the facets in \eqref{eq:facets} are the same three triangle
inequalities, so the polytope of \Cref{thm:B} coincides with the polytope of normalized entropies
$h/\log b$. The two observables are different, however. They agree on the mixtures
$\theta\delta_0+(1-\theta)\lambda$ of a point mass and Lebesgue measure $\lambda$, which have
normalized entropy $1-\theta$ and discrepancy $\theta$, but a periodic orbit of length at least two
has entropy zero and discrepancy less than one; in base $2$ the orbit of $1/3$ has $D_N\to2/3$.
Neither family of inequalities implies the other. Qualitative sufficiency in \Cref{thm:B} is also a
short consequence of Sigmund's generic-point theorem for systems with specification
\cite{Sigmund74} applied to the five vertex laws, together with weak continuity of $D$; we give a
formal proof of it and an arithmetic proof with a rate. We have not located \Cref{thm:B}, the
characterization in \Cref{thm:A}, or the scalar statements of \Cref{thm:C,thm:D} for extreme
discrepancy in the sources examined; this is a statement about those sources only. A comparison
with each source is kept in the ledger accompanying the source files (\Cref{sec:formal}).

\paragraph{Companion manuscript and evidence labels.}
The constructions of \Cref{thm:B,thm:D} use the positive state recurrence and the prime-power
orbit count of the companion manuscript \cite{CompanionC}, which in turn cites \Cref{thm:sharp}
from the present paper. Each numbered statement carries one label: \emph{Kernel-verified} (the
complete statement is a Lean declaration checked by the kernel, linked at a pinned commit),
\emph{Lean-supported paper proof} (named steps are Lean declarations), \emph{Paper proof}, or
\emph{External input}.

\section{Preliminaries}\label{sec:prelim}

For a finite sample $u=(u_0,\ldots,u_{N-1})$ in $\T$, identified with $[0,1)$, let $\mu_u$ be its
empirical probability measure. For a Borel probability measure $\mu$ on $\T$ put
$D(\mu)=\sup_{0\le a<c\le1}|\mu([a,c))-(c-a)|$ and $D(u)=D(\mu_u)$, so that
$D_N(x)=D((\fr{b^nx})_{n<N})$. A \emph{half-open arc} is the image in $\T$ of $[\theta,\theta+l)$ with
$0\le l\le1$; its length is $|I|=l$ and its \emph{excess} under $\mu$ is $\mu(I)-|I|$.

\begin{lemma}[Arcs]\label{lem:arcs}\evidence{Paper proof}
For every Borel probability measure $\mu$ on $\T$, $D(\mu)=\sup_I(\mu(I)-|I|)$, the supremum over
half-open arcs, and $|\mu(I)-|I||\le D(\mu)$ for every half-open arc. In particular $D(\mu)$ is
invariant under rotations of $\T$, and $D$ of a sample is unchanged when every point is rotated by
the same amount.
\end{lemma}

\begin{proof}
A non-wrapping arc is an interval $[a,c)$. A wrapping arc is the complement of an interval, and its
excess is the negative of that interval's excess; the full circle has excess zero. Hence the
supremum of $|\mu(I)-|I||$ over arcs is $D(\mu)$, and every negative excess $-e$ is matched by the
excess $e$ of the complementary arc.
\end{proof}

\begin{lemma}[Finite comparison]\label{lem:compare}\evidence{Paper proof}
Let $u,v$ be samples of length $N$, let $\epsilon\ge0$, and let
$E=\#\{n<N:\dist{u_n-v_n}>\epsilon\}$. Then $|D(u)-D(v)|\le E/N+2\epsilon$.
\end{lemma}

\begin{proof}
Let $I=[\theta,\theta+l)$. If $l>1-2\epsilon$ then $\mu_u(I)-l<2\epsilon$. Otherwise every $n$ with
$u_n\in I$ and $\dist{u_n-v_n}\le\epsilon$ has $v_n$ in the half-open arc $[\theta-\epsilon,
\theta+l+\epsilon+\eta)$ for every $\eta>0$, so $\mu_u(I)\le l+2\epsilon+\eta+D(v)+E/N$ by
\Cref{lem:arcs}. Let $\eta\to0$, take the supremum over $I$, and exchange $u$ and $v$.
\end{proof}

For radix orbits $u_n=\fr{b^nx}$ and $v_n=\fr{b^n(x+z)}$ the exceptional set depends only on $z$;
this case is the kernel-verified \lean{Phase18Synthetic}{583}{\ttt{discrepancy\_transfer}}. Rational
models are compared with real orbits through the following form.

\begin{lemma}[Model comparison]\label{lem:model}\evidence{Paper proof}
Let $x\in\R$, $N\ge1$, let $y_0,\ldots,y_{N-1}\in\T$ and let $\nu$ be a Borel probability measure on
$\T$. Suppose that $|\#\{n<N:y_n\in I\}-N\nu(I)|\le H$ for every half-open arc $I$, and that
$\dist{b^nx-y_n}\le\epsilon$ for all but $E$ indices $n<N$. Then
$|D_N(x)-D(\nu)|\le(H+E)/N+2\epsilon$.
\end{lemma}

\begin{proof}
By \Cref{lem:arcs}, $|D(y)-D(\nu)|\le\sup_I|\mu_y(I)-\nu(I)|\le H/N$. Apply \Cref{lem:compare}.
\end{proof}

We shall use two reference measures. If $0\le\theta\le1$ and $m\ge1$, let $U_m$ be the uniform
probability measure on $\{j/m:0\le j<m\}$. Then
\begin{equation}\label{eq:refdisc}
D\bigl(\theta\delta_0+(1-\theta)\lambda\bigr)=\theta,\qquad
D\bigl(\theta\delta_0+(1-\theta)U_m\bigr)=\theta+\frac{1-\theta}m .
\end{equation}
For the first, every arc has excess $\theta(\mathbf 1_{0\in I}-|I|)\in[-\theta,\theta]$, and short
arcs at $0$ approach $\theta$. For the second, a half-open arc containing $k\ge1$ grid points has
length greater than $(k-1)/m$, so its excess is less than $\theta+(1-\theta)k/m-(k-1)/m
=\theta+(1-\theta k)/m\le\theta+(1-\theta)/m$; arcs avoiding $0$ do no better, arcs containing no
grid point have nonpositive excess, and the arcs $[0,h)$ approach the bound as $h\downarrow0$. By
\Cref{lem:arcs} the excesses determine $D$.

\paragraph{Two tools from the companion manuscript.}
Fix an odd prime $p\nmid b$, $t=\ord_p(b)$ and $v=v_p(b^t-1)\ge1$.

\begin{lemma}[Prime-power orbits {\cite{CompanionC}}]\label{lem:period}\evidence{Lean-supported paper proof}
For $e\ge0$, $a$ prime to $p$, $j_0\ge0$, $L\ge1$ and every half-open arc $I$,
\[
\bigl|\#\{j_0\le j<j_0+L:\fr{ab^j/p^e}\in I\}-L|I|\bigr|\le p^vL/p^e+p^e .
\]
\end{lemma}

\begin{proof}
If $e\le v$ the left side is at most $L\le p^vL/p^e$. If $e=v+g$ with $g\ge1$, then over one period
of $tp^g$ consecutive indices the count of any interval $[a',c')$ differs from $tp^g(c'-a')$ by at most
$t$; this is the kernel-verified
\leanP{PrimePowerGrid}{313}{\ttt{primePower\_period\_count}}, applied to the numerator $ab^{j_0}$.
By \Cref{lem:arcs} the same holds for arcs. Writing $L=ftp^g+r$ with $0\le r<tp^g$ gives an error at
most $f t+r\le L/p^g+tp^g\le p^vL/p^e+p^e$, since $t<p$.
\end{proof}

\begin{lemma}[Positive state programming {\cite{CompanionC}}]\label{lem:program}\evidence{Kernel-verified}
Let $0=n_0<n_1<\cdots$ be integers and $u_0=0,u_1,u_2,\ldots$ rationals in $[0,1)$. Define $R_0=0$,
let $c_k\in(0,1]$ be the representative of $u_k-b^{n_k-n_{k-1}}R_{k-1}$ modulo one, and put
$R_k=b^{n_k-n_{k-1}}R_{k-1}+c_k$. Then $R_k=b^{n_k}\sum_{i\le k}c_ib^{-n_i}$ and $\fr{R_k}=u_k$ for
every $k$. If every $u_k$ has denominator dividing a power of $p$, so has every $c_k$.
\end{lemma}

In Lean the statement consists of
\lean{Phase18Synthetic}{383}{\ttt{programState\_fract}},
\lean{Phase18Synthetic}{438}{\ttt{programState\_eq\_prefix}},
\lean{Phase18Synthetic}{459}{\ttt{programmed\_prefix\_fract}},
\lean{Phase18Synthetic}{364}{\ttt{programCoeff\_mem}} and
\lean{Phase18Synthetic}{426}{\ttt{programCoeffs\_prime\_power}}; its two-line inductive proof is in
\cite{CompanionC}. More generally, $c_k$ has denominator dividing $\mathrm{lcm}(\den u_k,\den u_{k-1})$.
For $x=\sum_kc_kb^{-n_k}$ and $n_k\le n<n_{k+1}$ the rational part of $b^nx$ is $b^{n-n_k}R_k$, whose
fractional part is $\fr{b^{n-n_k}u_k}$. Since $0<c_i\le1$ and the positions are distinct,
\begin{equation}\label{eq:tail}
0<b^nx-b^{n-n_k}R_k\le\frac b{b-1}\,b^{-(n_{k+1}-n)} .
\end{equation}

\section{The sharp translation identity}\label{sec:translation}

\begin{theorem}[Sharp translation identity]\label{thm:sharp}\evidence{Kernel-verified}
For all real $x,z$ and every $N\ge1$,
\[
|D_N(x+z)-D_N(x)|\le1-D_N(z),
\]
with equality at $x=0$. Hence $\sup_x|D_N(x+z)-D_N(x)|=1-D_N(z)$, and the supremum is attained.
\end{theorem}

\begin{proof}
Let $J=[\gamma,\gamma+l)$ be a half-open arc with centre $c=\gamma+l/2$, and let $1-e$ be the
proportion of $n<N$ with $\fr{b^nz}\in J$. For those $n$,
$\dist{b^n(x+z)-(b^nx+c)}=\dist{b^nz-c}\le l/2$. The sample $(\fr{b^nx+c})_{n<N}$ has discrepancy
$D_N(x)$ by \Cref{lem:arcs}, so \Cref{lem:compare} with $\epsilon=l/2$ gives
$|D_N(x+z)-D_N(x)|\le e+l=1-(\mu_N^z(J)-|J|)$, where $\mu_N^z$ is the empirical measure of the orbit
of $z$. Taking the infimum over $J$ and using \Cref{lem:arcs} gives the inequality. The orbit of $0$
is constant, so $D_N(0)=1$ (test $[0,h)$ with $h\downarrow0$), and at $x=0$ the left side equals
$1-D_N(z)$. Lean declarations: \lean{SharpTranslation}{343}{\ttt{sharp\_translation}},
\lean{SharpTranslation}{357}{\ttt{radixDiscrepancy\_zero}},
\lean{SharpTranslation}{378}{\ttt{sharp\_translation\_zero}} and
\lean{SharpTranslation}{386}{\ttt{sharp\_translation\_isGreatest}}.
\end{proof}

The proof uses the orbit structure only through the pairing $b^n(x+z)=b^nx+b^nz$: for arbitrary
paired finite samples with $q_n\equiv p_n\pm r_n\pmod1$ it gives $D(q)\le D(p)+1-D(r)$
(\lean{SharpTranslation}{326}{\ttt{sampleDisc\_paired}}). For $\epsilon\ge0$ let $E_N(z,\epsilon)$
count the $n<N$ with $\dist{b^nz}>\epsilon$. The arcs $[-\epsilon,\epsilon+h)$ give
\begin{equation}\label{eq:deficit}
1-D_N(z)\le2\epsilon+E_N(z,\epsilon)/N .
\end{equation}

\begin{corollary}[Transfer]\label{cor:transfer}\evidence{Kernel-verified}
For all real $x,z$, $\epsilon\ge0$ and $N\ge1$, $|D_N(x+z)-D_N(x)|\le2\epsilon+E_N(z,\epsilon)/N$. If
$b^nz\in\Z$ for all $n\ge K$, then $|D_N(x+z)-D_N(x)|\le\min(1,K/N)$.
\end{corollary}

\begin{proof}
Combine \Cref{thm:sharp} with \eqref{eq:deficit}; for the second claim take $\epsilon=0$.
Lean declarations: \lean{Phase18Synthetic}{583}{\ttt{discrepancy\_transfer}} and
\lean{Phase18Synthetic}{642}{\ttt{discrepancy\_eventually\_int}}.
\end{proof}

\subsection{Universal preservers}

For real $w$ put $M_N(w)=N^{-1}\sum_{n<N}\dist{b^nw}$, and let $F_b=\{c\in\T:(b-1)c=0\}$ be the
set of the $b-1$ fixed points of $u\mapsto bu$ on $\T$.

\begin{theorem}[Universal numerical discrepancy preservers]\label{thm:preserve}\evidence{Lean-supported paper proof}
For real $z$ the following are equivalent.
\begin{enumerate}[label=\textup{(\roman*)},nosep]
\item $D_N(z)\to1$.
\item $\sup_{x\in\R}|D_N(x+z)-D_N(x)|\to0$.
\item $D_N(x+z)-D_N(x)\to0$ for every $x\in\R$.
\item For every $x$, $D_N(x+z)$ and $D_N(x)$ have the same lower limit and the same upper limit.
\item $\liminf_ND_N(z)=1$.
\item $z=c+w$ with $(b-1)c\in\Z$ and $M_N(w)\to0$.
\end{enumerate}
In that case the empirical measures of the orbit of $z$ converge weakly to the point mass at $c$
modulo one. In base $2$, condition \textup{(vi)} reads $M_N(z)\to0$.
\end{theorem}

\begin{proof}
The Lean-supported step is \Cref{thm:sharp}, which gives (i)$\Leftrightarrow$(ii).
(ii)$\Rightarrow$(iii)$\Rightarrow$(iv) are immediate; (iv)$\Rightarrow$(v) follows by taking $x=0$,
since $D_N(0)=1$; and (v)$\Rightarrow$(i) because $D_N\le1$.

(vi)$\Rightarrow$(i). Since $bc\equiv c\pmod1$, $\fr{b^nz}$ is the rotation by $c$ of $\fr{b^nw}$, so
$D_N(z)=D_N(w)$ by \Cref{lem:arcs}. Markov's inequality gives $E_N(w,\epsilon)\le NM_N(w)/\epsilon$,
and \eqref{eq:deficit} gives $1-D_N(w)\le2\epsilon+M_N(w)/\epsilon$ for every $\epsilon>0$.

(i)$\Rightarrow$(vi). Choose half-open arcs $J_N$ with $\mu_N^z(J_N)-|J_N|\ge D_N(z)-1/N$, and let
$c_N$ be the centre of $J_N$. Then $|J_N|\to0$ and $\mu_N^z(J_N)\to1$. Let $G_N$ be the set of $n<N-1$
for which $\fr{b^nz}$ and $\fr{b^{n+1}z}$ both lie in $J_N$; then $\#G_N=N-o(N)$. For $n\in G_N$,
\[
\dist{(b-1)c_N}\le\dist{b(b^nz-c_N)}+\dist{b^{n+1}z-c_N}\le(b+1)|J_N|,
\]
so the distance from $c_N$ to $F_b$ tends to zero. The points of $F_b$ are $1/(b-1)$ apart, so for
large $N$ there is a unique $f_N\in F_b$ nearest to $c_N$, and $J_N$ lies in a small neighbourhood of
$f_N$. If $f_{N+1}\ne f_N$ for infinitely many $N$, then for such $N$ the arcs $J_N$ and $J_{N+1}$
are disjoint, while $\mu_{N+1}^z(J_N)\ge\frac N{N+1}\mu_N^z(J_N)\to1$ and $\mu_{N+1}^z(J_{N+1})\to1$,
which is impossible. Hence $f_N=c$ for all large $N$, and for every $\rho>0$ the proportion of
$n<N$ with $\dist{b^nz-c}>\rho$ tends to zero. Lift $c$ to $\R$ with $(b-1)c\in\Z$ and put $w=z-c$.
Then $\dist{b^nw}=\dist{b^nz-c}$, so $M_N(w)\le\rho+o(1)$ for every $\rho>0$. The same estimate gives
the weak convergence of the empirical measures to $\delta_c$.
\end{proof}

\begin{remark}[Comparison with Rauzy's theorem]\label{rem:rauzy}
Let $\mathcal D(b)$ be the set of $z$ such that $x+z$ is normal to base $b$ for every normal $x$.
Rauzy proved that $\mathcal D(b)$ is the set of numbers of zero noise \cite{Rauzy76}; in the
formulation of Bergelson and Downarowicz these are the deterministic numbers, whose base-$b$ digit
sequences have only zero-entropy empirical limit measures \cite[Definitions~3.5 and~3.10,
Theorem~3.12]{BD25}. The set $K_b$ of \Cref{thm:preserve} is contained in $\mathcal D(b)$: if
$M_N(w)\to0$, then for each $l$ the length-$l$ digit blocks of $w$ are constant blocks $0^l$ or
$(b-1)^l$ at all but $o(N)$ positions, so every empirical limit measure of the digit sequence is
supported on the two fixed points of the shift; since $\mathcal D(b)$ is closed under addition and
contains the rationals, $c+w\in\mathcal D(b)$. The
inclusion is strict: in base $2$, $z=1/3$ is rational, hence in $\mathcal D(2)$, but its orbit
alternates between $1/3$ and $2/3$, so $D_N(1/3)\to2/3$ and $\sup_x|D_N(x+1/3)-D_N(x)|\to1/3$.
Preserving normality does not imply preserving the numerical discrepancy of nonnormal inputs.
\end{remark}

\subsection{Sparse translations}

\Cref{cor:transfer} turns sparseness of the digits of $z$ into a uniform rate.

\begin{theorem}[A quadratic sparse remainder]\label{thm:quadratic}\evidence{Lean-supported paper proof}
Let $Z=\sum_{k\ge1}2^{-k(k+1)}$. For every real $x$, every $N\ge1$ and every integer $w\ge0$,
\[
|D_N(x+Z;2)-D_N(x;2)|\le4\cdot2^{-w}+\frac{w\lfloor\sqrt{N+w}\rfloor}{N}.
\]
Hence $\sup_x|D_N(x+Z)-D_N(x)|=1-D_N(Z)=O(N^{-1/2}\log(2N))$. The number $Z$ is irrational, and
$x+Z$ is normal to base $2$ whenever $x$ is.
\end{theorem}

\begin{proof}
The displayed finite inequality is kernel-verified as
\lean{Phase18Synthetic}{837}{\ttt{quadratic\_sparse\_transfer}}. At a shift $n$ with no position
$k(k+1)$ in $(n,n+w]$, $2^nZ$ is within $2\cdot2^{-w}$ of an integer; at most $w$ shifts precede
each position, and at most $\lfloor\sqrt{M}\rfloor$ positions lie below $M$. Taking
$w=\lceil\log_2(N+1)\rceil$ gives the rate, and \Cref{thm:sharp} identifies the supremum. The binary
expansion of $Z$ has infinitely many ones and unbounded runs of zeros, so it is not eventually
periodic.
\end{proof}

\begin{proposition}[A factorial sparse remainder]\label{prop:factorial}\evidence{Lean-supported paper proof}
Let $Z_!=\sum_{k\ge2}b^{-k!}$. Then $Z_!$ is a Liouville number, and
$\sup_x|D_N(x+Z_!)-D_N(x)|=O(N^{-1}\log^2(2N))$.
\end{proposition}

\begin{proof}
The kernel-verified \lean{Phase18Synthetic}{811}{\ttt{sparse\_discrepancy\_transfer}} with unit
coefficients and $\epsilon=2b^{-w}$ gives the bound $4b^{-w}+w\,\#\{k:k!<N+w\}/N$; take
$w=\lceil\log_b(N+1)\rceil$ and note $\#\{k:k!<M\}\le1+\log_2M$. The partial sums
$P_k=\sum_{2\le i\le k}b^{-i!}$ have denominator $b^{k!}$ and
$0<Z_!-P_k\le\frac b{b-1}b^{-(k+1)!}$, which is smaller than $(b^{k!})^{-\mu}$ for every fixed $\mu$
and all large $k$.
\end{proof}

For arithmetic remainders \Cref{thm:sharp} reduces uniform translation rates to an estimate of the
single deficit $1-D_N(z)$. For example, the binary Lambert series $\sum_m\chi(m)/(2^m-1)$ of a
nonprincipal real character, and its multiplicative generalizations, have deficits of order
$(\log\log N)/(\log N)^\kappa$; this estimate and its translation consequence are proved in the
companion manuscript \cite{CompanionLambert}.

\section{The discrepancy distance}\label{sec:metric}

\begin{proposition}[Finite discrepancy distance]\label{prop:metric}\evidence{Kernel-verified}
For $N\ge1$ put $d_N(x,y)=1-D_N(x-y)$. Then $d_N$ is a translation-invariant pseudometric on $\R$,
\[
d_N(x,y)=\max_{w\in\R}|D_N(w+x-y)-D_N(w)|,
\]
and $|D_N(x)-D_N(y)|\le d_N(x,y)$. For $N\ge2$, $d_N(x,y)=0$ if and only if $(b-1)(x-y)\in\Z$.
Moreover $d_N(mx,my)\le|m|d_N(x,y)$ for every integer $m$.
\end{proposition}

\begin{proof}
The maximum formula is \Cref{thm:sharp}. Applying \Cref{thm:sharp} with origin $y-z$ and translation
$x-y$ gives the triangle inequality. Symmetry follows from \Cref{thm:sharp} at $x=-z$, which gives
$1-D_N(-z)\le1-D_N(z)$ and, by exchanging $z$ and $-z$, equality. If $N\ge2$ and $D_N(u)=1$, there
are half-open arcs of length tending to zero containing all $N$ points, so $\fr{u}=\fr{bu}$, that is
$(b-1)u\in\Z$; conversely such $u$ has a constant orbit. The integer Lipschitz bound follows from
the triangle inequality and translation invariance. Lean declarations:
\leanP{DiscrepancyGeometry}{55}{\ttt{discDist\_translate}},
\leanP{DiscrepancyGeometry}{68}{\ttt{discDist\_comm}},
\leanP{DiscrepancyGeometry}{78}{\ttt{discDist\_triangle}},
\leanP{DiscrepancyGeometry}{87}{\ttt{discDist\_isGreatest}},
\leanP{DiscrepancyGeometry}{93}{\ttt{radixDiscrepancy\_lipschitz}},
\leanP{DiscrepancyGeometry}{243}{\ttt{discDist\_eq\_zero\_iff}} and
\leanP{DiscrepancyGeometry}{346}{\ttt{discDist\_intMul\_le}}.
\end{proof}

The zero set and the local behaviour can be made explicit. Write $z=x-y$. For $N\ge2$,
\begin{equation}\label{eq:local}
\min\bigl(1/N,\dist{(b-1)z}\bigr)\le d_N(x,y)\le\frac{b^{N-1}-1}{b-1}\dist{(b-1)z}.
\end{equation}
Indeed, write $z=h+u$ with $(b-1)h\in\Z$ and $(b-1)|u|=\dist{(b-1)z}$. Since $bh\equiv h$, the samples
$h+b^nu$, $n<N$, lie on a real segment of length $(b^{N-1}-1)|u|$, which an arc of slightly larger
length covers; this gives the upper bound. An arc that misses a sample has excess at most $1-1/N$,
and an arc that contains the first two samples has length at least their circle distance
$\dist{(b-1)u}=\dist{(b-1)z}$; this gives the lower bound. Thus at each finite $N\ge2$ the quotient
of $(\R,d_N)$ is the circle $\R/(b-1)^{-1}\Z$ with its usual topology.

\subsection{Classical ingredients}

For a Borel probability measure $\mu$ on $\T$ let $Q_\mu(l)=\sup\{\mu(I):|I|=l\}$, the supremum over
half-open arcs; on the line this is L\'evy's concentration function \cite{Levy39}. Define the Ky Fan
distance of $\mu$ to constants by
\[
k(\mu)=\inf\{\epsilon\ge0:\ \mu(\{u:\dist{u-c}>\epsilon\})\le\epsilon\text{ for some }c\in\T\}.
\]
If $U$ is a random point with law $\mu$, then $k(\mu)$ is the infimum over constants $c$ of the Ky Fan
distance between $U$ and $c$ \cite{Fan43}, and, up to the convention at boundary points, the
infimum over $c$ of the Lévy--Prokhorov distance between $\mu$ and $\delta_c$
\cite{Prokhorov56}, \cite[Section~11.3]{Dudley02}.

\begin{proposition}[Concentration]\label{prop:concentration}\evidence{Paper proof}
For every Borel probability measure $\mu$ on $\T$,
\[
D(\mu)=\sup_{0\le l\le1}\bigl(Q_\mu(l)-l\bigr),\qquad
1-D(\mu)=\inf_{c\in\T,\ 0\le r\le1/2}\bigl(\mu(\{u:\dist{u-c}>r\})+2r\bigr),
\]
and $k(\mu)\le1-D(\mu)\le3k(\mu)$. Both constants are sharp.
\end{proposition}

\begin{proof}
The first formula is \Cref{lem:arcs}. A closed arc of length $2r<1$ is contained in half-open arcs of
length $2r+h$ for every $h>0$, and contains half-open arcs of length $2r$; hence the supremum of the
excess over half-open arcs equals the supremum of $\mu(B)-2r$ over closed balls
$B=\{\dist{u-c}\le r\}$, which is the second formula (for $r=1/2$ both sides are trivial).

Given $c$ and $r$ with $\mu(\dist{u-c}>r)=E$, the tolerance $\epsilon=\max(E,r)$ is admissible in the
definition of $k$, so $k(\mu)\le E+2r$; taking the infimum gives $k\le1-D$. Conversely the tolerance
$1/2$ is always admissible, and an admissible tolerance $\epsilon\le1/2$ at $c$ gives
$1-D(\mu)\le\epsilon+2\epsilon$. For Lebesgue measure $k=1/3$ and $1-D=1$. For
$(1-h)\delta_0+h\delta_{1/2}$ with $0<h<1/4$, a ball of radius less than $h$ cannot contain both atoms
and misses mass at least $h$, so $k=h$, while $c=0$, $r=0$ give $1-D\le h$; hence $1-D=k$.
\end{proof}

In this language \Cref{thm:sharp} says that the largest change of discrepancy caused by the
translation $x-y$ is the concentration deficit of the orbit of $x-y$ about its best centre. The
construction is a finite-sample identity, and it uses no independence: the paired orbits of $x$ and
$x+z$ are coupled by the arithmetic of $\R$.

\subsection{The asymptotic kernel and rational scalars}

Put $d^+(x,y)=\limsup_Nd_N(x,y)=1-\liminf_ND_N(x-y)$, and recall $K_b=\{z:D_N(z)\to1\}$.

\begin{proposition}[Asymptotic distance]\label{prop:asymptotic}\evidence{Lean-supported paper proof}
$d^+$ is a translation-invariant pseudometric, $d^+(mx,my)\le|m|d^+(x,y)$ for integers $m$, and
$d^+(x,y)=\max_w\limsup_N|D_N(w+x)-D_N(w+y)|$. Its kernel $\{z:d^+(z,0)=0\}$ is $K_b$, which is a
subgroup of $\R$ equal to $F_b+\{w:M_N(w)\to0\}$ (with $F_b$ lifted to $(b-1)^{-1}\Z$).
\end{proposition}

\begin{proof}
The pseudometric properties and the Lipschitz bound are
\leanP{DiscrepancyGeometry}{390}{\ttt{asymDist\_triangle}} and
\leanP{DiscrepancyGeometry}{419}{\ttt{asymDist\_intMul\_le}}, with the companion declarations at the
same commit. The maximum formula follows from \Cref{prop:metric} and the choice $w=-y$. The kernel
consists of the $z$ with $\liminf D_N(z)=1$, which is $K_b$ by \Cref{thm:preserve}, and a kernel of a
translation-invariant pseudometric is a subgroup.
\end{proof}

The kernel is dense in $\R$, since it contains every rational whose denominator divides a power of
$b$; the ordinary quotient topology on $\R/K_b$ is therefore indiscrete, while $d^+$ separates, for
instance, $0$ from every normal number $x/b^j$, however small.

\begin{lemma}\label{lem:rationalK}\evidence{Paper proof}
A rational number $r$ belongs to $K_b$ if and only if its reduced denominator divides $b^j(b-1)$ for
some $j\ge0$. Moreover $|D_N(z/b)-D_N(z)|\le1/N$ for all real $z$.
\end{lemma}

\begin{proof}
The orbit of $r$ is eventually periodic and its empirical measures converge to the uniform measure on
the terminal cycle. By \Cref{thm:preserve}, $r\in K_b$ exactly when this limit is a point mass at a
fixed point, that is, when $b^jr\equiv c\pmod 1$ for some $j$ and some $c$ with $(b-1)c\in\Z$; this
means $(b-1)b^jr\in\Z$. For the second claim, the samples of $z/b$ are $\fr{z/b}$ followed by the
first $N-1$ samples of $z$, so the two samples differ in one point and arc counts differ by at most
one.
\end{proof}

\begin{theorem}[Universal rational scalars]\label{thm:scalars}\evidence{Lean-supported paper proof}
Let $a/d$ be a reduced fraction with $d\ge1$ and $a\ne0$. Then $(a/d)K_b\subseteq K_b$ if and only if
every prime factor of $d$ divides $b$. Equivalently, the rational scalars preserving the kernel of
$d^+$ form the ring $\Z[1/b]$.
\end{theorem}

\begin{proof}
If every prime factor of $d$ divides $b$, then $d\mid b^j$ for some $j$ and $a/d=(ab^j/d)/b^j$. Integer
multiples preserve $K_b$ because it is a subgroup (\Cref{prop:asymptotic}), and division by $b$
preserves it by \Cref{lem:rationalK}. Conversely suppose that a prime $\ell\mid d$ does not divide $b$.
The fixed point $1/(b-1)$ lies in $K_b$. The reduced denominator of $a/(d(b-1))$ has $\ell$-adic
valuation $v_\ell(d)+v_\ell(b-1)>v_\ell(b^j(b-1))$ for every $j$, because $\ell\nmid a$; by
\Cref{lem:rationalK}, $(a/d)\cdot1/(b-1)\notin K_b$. The Lean-supported step is \Cref{thm:sharp}, through
\Cref{thm:preserve}.
\end{proof}

For an individual element of $K_b$ the question is arithmetic.

\begin{proposition}[Residue-density test]\label{prop:residue}\evidence{Paper proof}
Let $z\in K_b$, let $c\in\R$ with $(b-1)c\in\Z$ be as in \Cref{thm:preserve}\textup{(vi)}, and let
$k_n$ be an integer nearest to $b^nz-c$. For a reduced fraction $a/d$ with $d\ge1$, the following are
equivalent: $(a/d)z\in K_b$; $z/d\in K_b$; there is a residue $r$ modulo $d$ such that
$\#\{n<N:k_n\equiv r\pmod d\}=N-o(N)$.
\end{proposition}

\begin{proof}
If $z/d\in K_b$ then $(a/d)z=a(z/d)\in K_b$. Conversely, choose integers $\alpha,\beta$ with
$\alpha a+\beta d=1$; then $z/d=\alpha(a/d)z+\beta z\in K_b$ if $(a/d)z\in K_b$, since $K_b$ is a group.
Write $b^nz=c+k_n+e_n$ with $|e_n|\le1/2$; by \Cref{thm:preserve}, $|e_n|\to0$ outside a set of
density zero. Then $b^nz/d$ is within $|e_n|/d$ of $g_n=(c+k_n)/d$, which depends only on
$k_n\bmod d$, and the $d$ possible values of $g_n$ are distinct modulo one and $1/d$ apart. If one
residue has density one, the orbit of $z/d$ concentrates at one point and $D_N(z/d)\to1$. Conversely,
if $z/d\in K_b$ its orbit concentrates at one point by \Cref{thm:preserve}, so $g_n$ is constant modulo
one, and hence $k_n$ is constant modulo $d$, on a set of density one.
\end{proof}

The limiting point $(c+r)/d$ is automatically a fixed point. Since nonzero rational scalings preserve
normality \cite[Proposition~4.25]{BD25}, \Cref{prop:residue} gives elements of $K_b$ whose rational multiples
preserve normality under translation but change the numerical discrepancy of some input. For example,
in base $10$ let $S$ be the set of positive squares and $z=1/3+\sum_{a\in S}10^{-a}$, so that $c=1/3$.
Since $10\equiv1\pmod3$, $k_n$ is congruent modulo $3$ to $\lfloor\sqrt n\rfloor$ for all $n$ outside a set
of density zero, and $\lfloor\sqrt n\rfloor$ is equidistributed modulo $3$ in density; hence $z\in K_{10}$
but $z/3\notin K_{10}$.

\section{The addition spectrum}\label{sec:addition}

\begin{theorem}[Exact addition spectrum]\label{thm:addition}\evidence{Kernel-verified}
For every $N\ge1$ and all real $x,y$, the triple $(D_N(x),D_N(y),D_N(x+y))$ satisfies
\eqref{eq:facets}. Conversely, in every base $b\ge2$, for every $(\alpha,\beta,\gamma)\in[0,1]^3$
satisfying \eqref{eq:facets} there are real $x,y$ with $D_N(x)\to\alpha$, $D_N(y)\to\beta$ and
$D_N(x+y)\to\gamma$.
\end{theorem}

\begin{proof}
By \Cref{thm:sharp}, $|\gamma-\alpha|\le1-\beta$ and $|\gamma-\beta|\le1-\alpha$; together with
$\gamma\ge0$ these are \eqref{eq:facets}. The region $P$ that \eqref{eq:facets} defines in $[0,1]^3$ is
the convex hull of $(1,1,1)$, $(0,1,0)$, $(1,0,0)$, $(0,0,1)$ and $(0,0,0)$, and these are its extreme
points; explicit convex weights are given in \eqref{eq:weights} below. The converse is proved in Lean
by a block schedule adapted to complete periods; \Cref{thm:additionrate} gives a different proof with a
rate. Lean declarations: \leanP{AdditionSpectrum}{248}{\ttt{finite\_necessity}},
\leanP{AdditionSpectrum}{150}{\ttt{polytope\_eq\_convexHull}},
\leanP{AdditionSpectrum}{240}{\ttt{isExtremeIn\_polytope\_iff}},
\leanP{AdditionSpectrum}{1480}{\ttt{addSpectrum\_eq}} and
\leanP{AdditionSpectrum}{1486}{\ttt{exact\_addition\_spectrum}}. The formal proof of sufficiency
rests on the more general \leanP{AdditionSpectrum}{1470}{\ttt{mixtureRealizable}}: for integer
vectors $a^{(1)},\ldots,a^{(r)}\in\Z^d$ and convex weights $w_1,\ldots,w_r$ there are real
$x_1,\ldots,x_d$ such that every integer form $\sum_jc_jx_j$ has limiting discrepancy
$\sum\{w_i:c\cdot a^{(i)}=0\}$.
\end{proof}

In deficit coordinates $e=1-\delta$ the inequalities \eqref{eq:facets} read $e_1\le e_2+e_3$,
$e_2\le e_1+e_3$, $e_3\le e_1+e_2$; as explained in \Cref{sec:intro}, this is the polytope of the
normalized point-entropy inequalities of \cite[Proposition~4.9]{BD25}.

\subsection{A uniform cubic rate}

For $(\alpha,\beta,\gamma)$ satisfying \eqref{eq:facets} put $\sigma=\max(0,(\alpha+\beta+\gamma-1)/2)$
and
\begin{equation}\label{eq:weights}
w_1=\sigma,\quad w_2=\beta-\sigma,\quad w_3=\alpha-\sigma,\quad w_4=\gamma-\sigma,\quad
w_5=1-\alpha-\beta-\gamma+2\sigma .
\end{equation}
The inequalities \eqref{eq:facets} say that $\alpha+\beta+\gamma-1\le2\min(\alpha,\beta,\gamma)$, so
the five weights are nonnegative; they sum to one. With $u\in\T$, the five \emph{types} prescribe the
pairs of states $(0,0)$, $(u,0)$, $(0,u)$, $(u,-u)$, $(u,u)$ for $(x,y)$, hence the states
$(0,0,0)$, $(u,0,u)$, $(0,u,u)$, $(u,-u,0)$, $(u,u,2u)$ for $(x,y,x+y)$. A coordinate is \emph{retired}
(state $0$) for $x$ exactly in types $1$ and $3$, for $y$ in types $1$ and $2$, and for $x+y$ in types $1$
and $4$, with total weights $\alpha$, $\beta$ and $\gamma$.

\begin{lemma}[Scheduler]\label{lem:scheduler}\evidence{Paper proof}
Let $w_1,\ldots,w_r\ge0$ sum to one. For $k\ge1$ let $C_i(k-1)$ be the number of $i$ among the types
of blocks $1,\ldots,k-1$, and give block $k$ a type maximizing $kw_i-C_i(k-1)$. Then
$-(r-1)\le C_i(k)-kw_i\le1$ for all $i$ and $k$, and a type of weight zero is never chosen. The same
holds if the maximization uses rationals $\widehat w_i(k)$ with $|\widehat w_i(k)-w_i|\le1/(4rk)$.
\end{lemma}

\begin{proof}
The deficits $kw_i-C_i(k-1)$ sum to one, so the largest is at least $1/r$; with the approximate
weights the chosen deficit is at least $1/r-2k/(4rk)>0$. After the choice the chosen overcount is
less than one, the others decrease, and all overcounts sum to zero, which gives the lower bound by
induction. A type of weight zero has deficit $-C_i\le0$.
\end{proof}

\begin{theorem}[Uniform rate]\label{thm:additionrate}\evidence{Lean-supported paper proof}
Let $p\nmid b$ be an odd prime and $t=\ord_p(b)$. There is $C_{b,p}$ such that for every
$(\alpha,\beta,\gamma)$ satisfying \eqref{eq:facets} there are rationals $c_k,d_k\in(0,1]$ with
$p$-power denominators for which $x=\sum_kc_kb^{-tk^3}$ and $y=\sum_kd_kb^{-tk^3}$ satisfy, for all
$N\ge1$,
\[
\max\bigl(|D_N(x)-\alpha|,\ |D_N(y)-\beta|,\ |D_N(x+y)-\gamma|\bigr)\le C_{b,p}N^{-1/3}.
\]
The empirical measures of the three orbits converge weakly to $\alpha\delta_0+(1-\alpha)\lambda$,
$\beta\delta_0+(1-\beta)\lambda$ and $\gamma\delta_0+(1-\gamma)\lambda$.
\end{theorem}

\begin{proof}
Let $n_k=tk^3$, $h_k=n_{k+1}-n_k=t(3k^2+3k+1)\le7tk^2$, $q_k=p^{1+\lfloor\log_pk\rfloor}$ (so
$k<q_k\le pk$) and $v=v_p(b^t-1)$. Block $k$ consists of the integers in $[n_k,n_{k+1})$, with
$n_0=0$; block $0$ has type $1$, and blocks $k\ge1$ receive types by \Cref{lem:scheduler} with the
weights \eqref{eq:weights}, using $u=1/q_k$ in block $k$. Define the coefficient sequences of $x$ and $y$
by \Cref{lem:program} from their state sequences. Then the rational part of $b^{n_k}(x+y)$ has
fractional part equal to the sum of the two states modulo one, and $0<c_k+d_k\le2$. The nonzero states
$\pm u$ and $2u$ have exact denominator $q_k$ because $p$ is odd.

Fix a coordinate $z\in\{x,y,x+y\}$ with target $\delta$, and let $tK^3\le N<t(K+1)^3$ with $K\ge1$. Let
$y_n$ be the fractional part of the rational model of $b^nz$ and $Z$ the number of retired indices below
$N$. The retired indicator of $z$ on blocks $1,2,\ldots$ is the sum of the indicators of two types, so by
\Cref{lem:scheduler} its prefix sums differ from $k\delta$ by at most $8$. Summation by parts with the
increasing weights $h_k$ gives
\[
|Z-\delta N|\le t+17h_K .
\]
Active blocks are orbit segments of $a/q_k$ with $a\in\{1,-1,2\}$, so by \Cref{lem:period} the active
indices meet every half-open arc $I$ in $(N-Z)|I|+\theta$ points with
$|\theta|\le H:=\sum_{k\le K}(p^vh_k/q_k+q_k)\le(7tp^v+p)K(K+1)/2$. Retired models are $0$. Hence, with
$\nu=(Z/N)\delta_0+(1-Z/N)\lambda$, the model counts differ from $N\nu(I)$ by at most $H$ on every arc,
and $D(\nu)=Z/N$ by \eqref{eq:refdisc}. By \eqref{eq:tail}, applied to $x$ and $y$, the real point
$b^nz$ is within $\frac{2b}{b-1}b^{-W}$ of $y_n$ unless $n$ is one of the $W$ indices preceding an insertion.
With $W=\lceil\log_b(\frac{2b}{b-1}(K+1))\rceil$, \Cref{lem:model} gives
\[
|D_N(z)-\delta|\le\frac{(7tp^v+p)K(K+1)/2+(K+1)W+t+17h_K}{N}+\frac2{K+1},
\]
which is $O_{b,p}(K^{-1})=O_{b,p}(N^{-1/3})$, uniformly in the triple. For $N<t$ use $D_N\le1$. For weak
convergence approximate a continuous test function uniformly by step functions on a fixed partition
into arcs, and use the same three estimates. The Lean-supported step is \Cref{lem:program}.
\end{proof}

The exponent $1/3$ comes from balancing $1/K$, $K^2/N$ and $q_k^{-1}$ on blocks of one type; we do not
claim that it is optimal.

\subsection{Computability and algebraic independence}

\begin{theorem}[Effective independent witnesses]\label{thm:independent}\evidence{Lean-supported paper proof}
In \Cref{thm:additionrate} the witnesses $x,y$ can be chosen algebraically independent over $\Q$, with
error at most $C'_{b,p}N^{-1/3}$ uniformly in the triple. In particular $x$, $y$ and $x+y$ are
transcendental at every point of the polytope, including its vertices. There is one algorithm that,
given Cauchy names of a computable triple satisfying \eqref{eq:facets}, computes such $x$ and $y$.
\end{theorem}

\begin{proof}
\emph{Computability.} The weights \eqref{eq:weights} are computable from names of the triple, so at
step $k$ one obtains rationals within $1/(20k)$ of them and applies the approximate rule of
\Cref{lem:scheduler}; the bounds of \Cref{thm:additionrate} hold verbatim. The coefficients are
rationals computed from the states, and by \eqref{eq:tail} the tail after insertion $k$ is at most
$\frac b{b-1}b^{-n_{k+1}}$, which is a computable modulus of convergence.

\emph{Overrides.} Fix the nominal type sequence. At blocks $k$ that are powers of two we allow arbitrary
pairs of states in $(q_k^{-1}\Z/\Z)^2$. Each overridden block is charged entirely as model count error;
through $N$ the extra charge is at most $\sum_{2^j\le K}h_{2^j}\le7t\sum_{2^j\le K}4^j\le\frac{28t}3K^2$.
By \Cref{lem:program} the next block recovers its programmed state exactly, and all coefficients stay in
$(0,1]$, so \eqref{eq:tail} and the bounds of \Cref{thm:additionrate} hold with a larger constant.

\emph{Avoiding relations.} Enumerate the nonzero polynomials $P_1,P_2,\ldots\in\Z[X,Y]$. Suppose that the
coefficients through index $J$ are fixed, and let $P=P_i$ have degrees $d_X,d_Y$ in $X,Y$. Choose a power
of two $k>J$ with $q_k>\max(d_X,d_Y)$ and $\frac b{b-1}b^{-h_k}<1/(2q_k)$, fill the blocks between $J$
and $k$ nominally, and consider in block $k$ the states $(a/q_k,a'/q_k)$ with $0\le a\le d_X$ and
$0\le a'\le d_Y$, each continued nominally forever. The resulting first coordinates $x_a$ satisfy
$b^{n_k}x_a\in a/q_k+(0,1/(2q_k))$ modulo one, so they are distinct, and likewise the $y_{a'}$. A nonzero
polynomial of these degrees cannot vanish on a grid of $(d_X+1)(d_Y+1)$ points with distinct
coordinates, so $P(x_a,y_{a'})\ne0$ for some candidate. To find one, extend all candidates in parallel:
every continuation of a prefix $S_L$ through index $L$, whatever the later choices, lies in
$[S_L,S_L+\frac b{b-1}b^{-n_{L+1}}]$, because all coefficients are in $(0,1]$. Interval evaluation of
$P$ on these boxes eventually certifies $P\ne0$ for some candidate; fix it, together with the nominal
choices through the certifying depth, and proceed to $P_{i+1}$. The certificate survives every later
choice, so $P_i(x,y)\ne0$ for every $i$. The Lean-supported step is \Cref{lem:program}.
\end{proof}

The triple $(x,y,x+y)$ satisfies the linear relation $X+Y-Z=0$; the theorem concerns the two free
inputs.

\subsection{Several summands and a six-form inequality}

For $d$ summands it is convenient to use the deficits $e=1-\delta$.

\begin{theorem}[$d$ summands]\label{thm:dsummands}\evidence{Lean-supported paper proof}
Let $d\ge1$ and $e_1,\ldots,e_{d+1}\in[0,1]$. There are real $x_1,\ldots,x_d$ such that
$1-D_N(x_i)\to e_i$ for $i\le d$ and $1-D_N(x_1+\cdots+x_d)\to e_{d+1}$ if and only if
$e_i\le\sum_{j\ne i}e_j$ for every $i\le d+1$.
\end{theorem}

\begin{proof}
Write $\epsilon_N(z)=1-D_N(z)$. By \Cref{thm:sharp}, $\epsilon_N(u+v)\le\epsilon_N(u)+\epsilon_N(v)$ and
$\epsilon_N(u+v)\ge\epsilon_N(u)-\epsilon_N(v)$. Iterating the first inequality bounds the deficit of the
sum; adding the other summands to $x_i$ one at a time in the second gives
$\epsilon_N(x_i)\le\epsilon_N(x_1+\cdots+x_d)+\sum_{j\ne i,\,j\le d}\epsilon_N(x_j)$. This proves necessity.

The set $P_d$ defined by the inequalities is the convex hull of $0$ and the $0$--$1$ vectors with at
least two ones. These vectors lie in $P_d$, which is convex. If an extreme point of $P_d$ satisfies some
$e_i=\sum_{j\ne i}e_j$, then it is the convex combination
$(1-e_i)\cdot0+\sum_{j\ne i}e_j(\mathbf1_i+\mathbf1_j)$ of listed vectors, hence one of them. Otherwise all
these inequalities are strict; a coordinate strictly between $0$ and $1$ could then be perturbed in both
directions, so the point is a vertex of the cube, and a vertex with a single one violates its own
inequality.

A listed vector with support $A\subseteq\{1,\ldots,d+1\}$, $|A|=r\ge2$, corresponds to an integer
direction $a\in\Z^d$ for which the entries of $(a_1,\ldots,a_d,a_1+\cdots+a_d)$ are nonzero exactly on
$A$: give $r-1$ of the selected positions the value $1$ and the last the value $-(r-1)$, so that the
selected values sum to zero, negate the value in position $d+1$ if it is selected, and read off $a$ from
the first $d$ positions. For convex weights $w_r$ on such directions (the zero vector corresponding to
$0\in P_d$), the kernel-verified \leanP{AdditionSpectrum}{1470}{\ttt{mixtureRealizable}} gives reals
$x_1,\ldots,x_d$ such that every integer form $\sum_jc_jx_j$ has limiting discrepancy
$\sum\{w_r:c\cdot a^{(r)}=0\}$. Applied to the coordinate forms and to $(1,\ldots,1)$, this gives in
position $i$ the deficit $\sum\{w_r:\text{entry }i\text{ of }(a^{(r)},\textstyle\sum_ja^{(r)}_j)\ne0\}$,
which is the prescribed convex combination.
\end{proof}

Testing triples separately does not describe the joint behaviour of several forms.

\begin{proposition}[Six forms]\label{prop:sixforms}\evidence{Paper proof}
For all real $x,y,z$ and $N\ge1$,
\[
D_N(x+y)+D_N(x+z)+D_N(y+z)\le2+2\min\bigl(D_N(x),D_N(y),D_N(z)\bigr).
\]
In particular the three pairwise sums of three normal numbers cannot all have discrepancy tending to
one, although each triple $(0,0,1)$ is admissible in \Cref{thm:addition}.
\end{proposition}

\begin{proof}
For an integer $m\ne0$, the preimage of a half-open arc under $u\mapsto mu$ consists of $|m|$ half-open
arcs of total length equal to its length, so $D_N(mx)\le|m|D_N(x)$. Put $u=x+y$, $v=x+z$, $w=y+z$. By
\eqref{eq:facets}, $D_N(u+v)\ge D_N(u)+D_N(v)-1$, and since $2x=(u+v)+(-w)$ and $D_N(-w)=D_N(w)$
(\Cref{prop:metric}), $D_N(2x)\ge D_N(u+v)+D_N(w)-1$. Hence
$D_N(u)+D_N(v)+D_N(w)\le2+D_N(2x)\le2+2D_N(x)$. Permute $x,y,z$.
\end{proof}

In base $3$ the inequality is attained in the limit by $x=y=z$ generic for
$\frac12(\delta_0+\delta_{1/2})$, an invariant measure since both points are fixed. The joint spectra of
finite families of rational linear forms are therefore not determined by their triples. Arithmetic
realizations of Haar measures on rational subtori, a computable polyhedral approximation of these joint
spectra, and joint box-discrepancy rates for full-dimensional targets are not treated here.

\section{Rational operations}\label{sec:operations}

\begin{proposition}[Finite dilation bounds]\label{prop:dilation}\evidence{Lean-supported paper proof}
For every integer $m\ne0$, real $x$ and $N\ge1$,
\[
\max\bigl(0,1-|m|(1-D_N(x))\bigr)\le D_N(mx)\le\min\bigl(1,|m|D_N(x)\bigr).
\]
\end{proposition}

\begin{proof}
The lower bound is $d_N(mx,0)\le|m|d_N(x,0)$, the kernel-verified
\leanP{DiscrepancyGeometry}{346}{\ttt{discDist\_intMul\_le}}. The upper bound was proved in
\Cref{prop:sixforms}.
\end{proof}

At $m$ prime to $b$ both bounds are attained asymptotically, as the next theorem shows; at $m=b^j$ one
has instead $|D_N(b^jx)-D_N(x)|\le j/N$, so the relevant factor is the part of $m$ prime to $b$.

\begin{theorem}[Boundary spectrum of dilation]\label{thm:boundary}\evidence{Lean-supported paper proof}
Let $m\ne0$ be an integer and $|m|=m_0m_b$ with $\gcd(m_0,b)=1$ and every prime factor of $m_b$ dividing
$b$.
\begin{enumerate}[label=\textup{(\alph*)},nosep]
\item If $D_N(mx)\to1$, then $\liminf_ND_N(x)\ge1/m_0$.
\item Let $p\nmid b$ be an odd prime, $t=\ord_p(b)$, and $\delta\in[1/m_0,1]$. There are rationals
$c_k\in(0,1]$ with denominators dividing $m_0$ such that $x=\sum_kc_kb^{-tk^3}$ satisfies
$|D_N(x)-\delta|+|1-D_N(mx)|\le C_{b,p,m}N^{-1/3}$ for all $N\ge1$.
\end{enumerate}
\end{theorem}

\begin{proof}
(a) By \Cref{prop:asymptotic,thm:scalars}, $m_0x=\pm(mx)/m_b\in K_b$, since $1/m_b\in\Z[1/b]$. By
\Cref{thm:preserve} there is $c$ with $(b-1)c\in\Z$ such that, for every $\rho>0$, all but $o(N)$ of the
$n<N$ satisfy $\dist{m_0b^nx-c}\le\rho$. For those $n$, $\fr{b^nx}$ is within $\rho/m_0$ of the set of
$m_0$ equally spaced solutions of $m_0u\equiv c\pmod1$. An arc of length $(1-2\rho)/m_0$ inside a gap of
this set therefore has empirical mass $o(1)$, so $\liminf D_N(x)\ge(1-2\rho)/m_0$ for every $\rho>0$.

(b) The grid $G=\{j/m_0\}\subset\T$ is permuted by $u\mapsto bu$, because $\gcd(b,m_0)=1$; let
$\Gamma_1,\ldots,\Gamma_s$ be its cycles, of lengths $\ell_1,\ldots,\ell_s$. Put
$\theta=(\delta-1/m_0)/(1-1/m_0)\in[0,1]$ (any $\theta$ if $m_0=1$), so that by \eqref{eq:refdisc} the
measure $\nu_\theta=\theta\delta_0+(1-\theta)U_{m_0}$ has $D(\nu_\theta)=\delta$. Use types $0$ (state $0$)
and $1,\ldots,s$ (state a chosen point of $\Gamma_i$), with weights $\theta$ and $(1-\theta)\ell_i/m_0$,
schedule them by \Cref{lem:scheduler} on the blocks of $n_k=tk^3$, and define $x$ by \Cref{lem:program}.
The coefficients have denominators dividing $m_0$. In a block of type $i\ge1$ the model runs through
$\Gamma_i$, so its counts on every arc differ from the block length times the uniform measure on
$\Gamma_i$ by at most $\ell_i\le m_0$. Summation by parts as in \Cref{thm:additionrate} shows that the
model counts differ from $N\nu_\theta(I)$ by $O_{b,m}(K^2)$, and \Cref{lem:model} with \eqref{eq:tail}
gives $|D_N(x)-\delta|=O(N^{-1/3})$. Every state is annihilated by $m$, so the model of $mx$ is $0$
throughout, and the real points of $mx$ are within $|m|\frac b{b-1}b^{-W}$ of $0$ off $(K+1)W$ indices;
\Cref{lem:model} with $\nu=\delta_0$ gives $1-D_N(mx)=O(N^{-1/3})$. The Lean-supported step is
\Cref{lem:program}.
\end{proof}

The kernel $K_b$ is a group, preserved by $\Z[1/b]$, but it is not closed under multiplication, even
within its mean-zero part $\{w:M_N(w)\to0\}$, whose rational elements form the ring $\Z[1/b]$.

\begin{proposition}[A product of sparse preservers]\label{prop:product}\evidence{Paper proof}
Let $A$ be the set of nonnegative integers whose binary digits vanish in all odd positions, let $r\ge2$,
and put $U=\sum_{a\in A}b^{-ra}$ and $V=\sum_{a\in A}b^{-2ra}$. Then $U$ and $V$ are irrational,
$M_N(U)\to0$ and $M_N(V)\to0$, and $\sup_x|D_N(x+U)-D_N(x)|+\sup_x|D_N(x+V)-D_N(x)|=O(\log(2N)/\sqrt N)$. Their product
is $UV=b^r/(b^r-1)$, which is not in $K_b$.
\end{proposition}

\begin{proof}
Every nonnegative integer is uniquely $a+2a'$ with $a,a'\in A$, by splitting its binary digits into even
and odd positions; absolute convergence gives $UV=\sum_{n\ge0}b^{-rn}$. Below $4^j$ there are $2^j$
elements of $A$, so the positions of the nonzero digits of $U-1$ and $V-1$ number $O(\sqrt X)$ below
$X$.
At a shift with no position in $(n,n+w]$ the remainder is within $2b^{-w}$ of an integer; this gives
$M_N\to0$, and \Cref{cor:transfer} with $w=\lceil\log_bN\rceil$ gives the rate. The digits contain
infinitely many ones and unbounded runs of zeros, so the numbers are irrational. Finally
$\fr{UV}=1/(b^r-1)$, whose reduced denominator is prime to $b$ and greater than $b-1$, so it does not
divide $b^j(b-1)$; by \Cref{lem:rationalK}, $UV\notin K_b$.
\end{proof}

For $b\ge3$ the rational elements of $K_b$ are already not closed under multiplication, since
$(b-1)^{-2}\notin K_b$; the content of \Cref{prop:product} lies in base $2$ and in the mean-zero
subgroup. Bergelson and Downarowicz construct, in base $2$ and with signed sparse digits, two numbers
whose digit changes or ones have density zero and whose product is normal \cite[Section~8]{BD25}; both
factors then have $M_N\to0$, so their example is a stronger statement of the same kind. The example
here uses positive digits, gives an explicit rational product and holds in every base.

\begin{remark}[Scope of the identities]\label{rem:scope}
\Cref{thm:sharp} uses rotation invariance of the extreme discrepancy and fails for its anchored and
multidimensional analogues. For the star discrepancy, with intervals $[0,s)$, take $b=2$, $N=3$,
$x=1/4$ and $z=-1/256$: then $D_N^*(x)=1/2$, $D_N^*(x+z)=61/192$ and $D_N^*(z)=63/64$, so the change
$1/6+1/64$ exceeds $1-D_N^*(z)=1/64$. For the extreme box discrepancy in $[0,1)^2$, take $b=2$, $N=2$,
$x=(1/3,0)$ and $z=(0,1/3)$: each has box discrepancy one, by thin strips, while $x+z$ gives the points
$(1/3,1/3)$ and $(2/3,2/3)$, whose box discrepancy is $8/9$. These values were computed exactly.
\end{remark}

\section{Further questions}\label{sec:questions}

\begin{question}
Is the exponent $1/3$ in \Cref{thm:additionrate,thm:boundary} optimal among positive rational
expansions on a fixed schedule? Can a witness achieve $o(N^{-1/3})$ for a triple in the interior of the
polytope?
\end{question}

\begin{question}
Describe the joint limiting spectrum of $(D_N(m\cdot x))_{m\in M}$ for a finite set $M$ of integer
vectors. By \Cref{prop:sixforms} it is not determined by its two- and three-term projections, and it
depends on $b$ through the invariant measures of $u\mapsto bu$ on the torus.
\end{question}

\begin{question}
Which real $\lambda$ satisfy $\lambda K_b\subseteq K_b$? \Cref{thm:scalars} answers this for rational
$\lambda$ only.
\end{question}

\section{Formal verification and data availability}\label{sec:formal}

The Lean~4 \cite{Lean4} declarations cited above, built with Mathlib \cite{Mathlib}, are in the
repository \url{https://github.com/eldernos/MathematicsResearch}. Links to the modules
\ttt{SharpTranslation} and \ttt{Phase18Synthetic} are pinned to the commit
\begin{center}\ttt{\leancommit},\end{center}
and links to \ttt{AdditionSpectrum}, \ttt{DiscrepancyGeometry} and \ttt{PrimePowerGrid} to the later
commit
\begin{center}\ttt{\phaseXIXcommit}.\end{center}
At both commits the full build completed, the axiom audit accepted only \ttt{propext},
\ttt{Classical.choice} and \ttt{Quot.sound}, and an independent kernel replay covered all modules. The
definition of $D_N$ in Lean (\ttt{radixDiscrepancy}) is the supremum over $0\le a<c\le1$ used here.

The kernel-verified statements are
\Cref{lem:program,thm:sharp,cor:transfer,prop:metric,thm:addition}, together with the finite
inequality in \Cref{thm:quadratic} and the lower bound in \Cref{prop:dilation}. The rates in
\Cref{thm:additionrate,thm:boundary}, the computability and independence in \Cref{thm:independent},
the preserver classification, the scalar results, the convex geometry in \Cref{thm:dsummands} and
\Cref{prop:concentration,prop:sixforms,prop:product} have written proofs. Exact rational computations
exercising the addition spectrum and its cubic witnesses are in \ttt{scripts/phase19\_addition.py} at
the second commit; they are finite checks, not proofs of the asymptotic statements. The theorem
verification map, the novelty ledger and the reproduction commands are in the files
\ttt{VERIFICATION.md} and \ttt{LEDGER.md} next to the source of this paper
(directory \ttt{papers/10-discrepancy-arithmetic} of the repository).

\section{Use of AI tools}

This manuscript was produced in an AI-assisted research project. Large language models proposed
conjectures and proof strategies, wrote the proofs and the text, implemented the exact computations
and Lean~4 formalizations, searched the literature and reviewed drafts across several model providers.
The repository owner directed the project and is responsible for its release. The Lean kernel checked
the declarations listed in the verification map. The written proofs, literature comparisons and
priority statements have been checked only by AI models from different providers. They have not been
checked by an independent human expert. Cross-model review is not peer review.

\bibliographystyle{plainurl}
\bibliography{references}
\end{document}
